% III. Програмна реализация.
\section{Програмна реализация}
\label{sec:implementation}

Изходният текст е разделен по отговорност: \code{include/stratus/} съдържа логиката в заглавни
файлове, \code{src/} реализациите и главните функции, \code{tests/} модулните тестове,
\code{infra/} описанието на инфраструктурата. Указател към файловете е даден в приложението.
Проектът се изгражда с нула задължителни външни зависимости: достатъчни са компилатор на C++20 и
CMake. Двете места, където обичайно се появява зависимост, рамката за тестове и рамката за
измерване, са заменени с малки собствени реализации.

\textbf{Изчислително ядро и работник.} Ядрото е свободно от заделяне на памет и връща контролна
сума заедно с точния брой изпълнени вътрешни итерации. Контролната сума позволява проверка на
отговора и не позволява на оптимизатора да премахне цикъл, чийто резултат никой не чете.
Работникът излага \code{/healthz}, \code{/work?size=\&iter=} и \code{/metrics}. Параметър извън
допустимия обхват води до код 400, а не до подмяна с най-близката стойност: тихата подмяна би
направила измерване с отхвърлен параметър да изглежда успешно.

HTTP сървърът използва пул с фиксиран размер, а не нишка на връзка. Изборът е съществен: при нишка
на връзка едновременността е тази, която клиентът е избрал, процесорът се презаписва, а сигналът
за натоварване престава да означава каквото и да било. При фиксиран пул капацитетът на една
реплика е точно броят нишки в пула.

\textbf{Преносимост.} Единствената платформено зависима част е \code{src/net.cpp}, където се
поемат разликите между Winsock и сокетите на BSD. Всичко над този файл е преносим C++20 без нито
един условен блок за платформа.

\textbf{Посредник и обединяване на метрики.} Посредникът изпълнява две задачи, защото и на двете
им трябва едно и също нещо: текущият списък с живи реплики. Той се получава чрез разрешаване на
името на услугата, тъй като вграденият DNS връща по един адрес на работеща реплика. Няма регистър
на услугите и няма презареждане на конфигурация. При обединяването етикетът с името на репликата
се премахва и сериите с еднакъв ключ се сумират, което превръща $N$ серии в една за набора.

\textbf{Контролер.} Самото решение е чиста функция и се проверява от пет модулни теста без
работещ стек. Контролен цикъл, наблюдаем само чрез гледане на живо разгърнат клъстер, не може нито
да бъде обсъждан, нито проверяван. Едно място заслужава коментар, защото първата му версия беше
сгрешена: общата метрика е сума от броячите на репликите, така че премахването на реплика кара
сумата да \emph{спадне}, а първата версия прочете отрицателната разлика като отрицателно
натоварване и намали набора още повече. Правилното поведение е това на всяка система за
наблюдение: отрицателна разлика в брояч е нулиране и интервалът се пропуска.

\textbf{Генератор.} При \code{--rate 0} контурът е затворен: всяка нишка изпраща следващата заявка
веднага след връщането на предишната, тоест измерва се капацитет. При зададена скорост контурът е
отворен: моментите на изпращане са предварително определени, така че забавена услуга натрупва
време на чакане в отчетеното закъснение, вместо клиентът тихо да намали скоростта си. Персентилите
се изчисляват върху пълната извадка, а не от кофите на хистограма, защото границата на кофа,
отчетена като персентил, кара две различни системи да изглеждат еднакви.

\textbf{Модел на разхода.} Калкулаторът приема всяка цена като задължителен вход и отказва да
работи без нея. В изходния текст няма нито една публикувана цена: цените се различават по регион,
сметка и месец, така че цена, вписана в програма, е невярна още на следващия ден.

\textbf{Тестове и контейнер.} Тестовете са 26 и покриват детерминизма и цената на ядрото, разбора
на заявки и параметри, представянето и разчитането на метриките, обединяването между реплики,
петте случая на политиката, модела на разхода и статистиката за закъснението. Гетъри не се тестват.
Всеки тест е отделен запис в CTest, така че при провал се съобщава кой тест е паднал. Образът се
изгражда на два етапа: първият свързва изпълнимите файлове статично срещу musl, вторият е
\code{scratch} и съдържа само тях; измереният размер е \textbf{8{,}09 MB}. Тестовете се изпълняват
в етапа на изграждане: образ, който не минава тестовете си, не бива да съществува.

\textbf{Описание на инфраструктурата.} Кореновата конфигурация избира модул според стойността на
променлива за доставчик и подава към него един и същ обект; изразите с \code{count} в
\code{infra/main.tf} са единственото място, където доставчик се назовава извън модула си.
Количествената мярка за преносимост, преброена в редове без празни и без коментари, е: обща част
82 реда, модул за AWS 297 реда, модул за Azure 178 реда. Съотношението не е в полза на
твърдението, че едно описание покрива двамата доставчици: общата част е кратка, защото описва
\emph{услугата}, а мрежата, платформата, хранилището и идентичността се пишат два пъти.
\textbf{Двата модула са написани и форматирани, но не са прилагани.} В хранилището не се записват
идентификатори на сметки или ключове.
